\section{Section A: Conceptual Understanding}

\subsection{Question 1: Regularistion and Overfitting}

\begin{enumerate}[label=\textbf{(\alph*)}, leftmargin=*, align=left]
    \item Overfitting occurs when a model learns the characteristics of the training set too closely resulting in the model not learning underlying general patterns. As a result, 
          the model fails to generalise when exposed to an unseen (test) dataset. 
    
          For a sentiment classifier trained on the Polarity Dataset v2.0, two observable signs of overfitting the training data would be a poor F1-score and a low AUC-ROC score for a test set.
          A low F1-score indicates poor Precision and poor Recall which implies that the model has a high number of false positives and false negatives. Another sign would be a low AUC-ROC score.
          This indicates how close a model is to being perfect and indicates the ratio of true positives to false positives. An AUC-ROC score of close to 0.5 would indicate a model that is not that
          much better than a random model. This can be as a result of memorising specific words, phrases or patterns from the training set reviews.

          Overfitting is especially an issue in high-dimensional text feature spaces. As the dimension increases, data becomes more spread out which can make it harder for the model to 
          learn general patterns resulting in the model memorising instead. In addition, there can be instances in which the model find rare words that may appear predictive in training
          but are actually noise due to each word being treated as a feature. Lastly, many features (words) may not even be present in some reviews. The model can overemphasise these words
          also resulting in the model learning noise and hence overfitting.

    \newpage
    
    \item   Table \ref{reg_comp} illustrates the comparision between L1 and L2 regularisation:
    
            \begin{table}[htbp]
            \caption{L1 vs L2 Regularisation}
            \centering
            \begin{tabularx}{\columnwidth}{p{5.2cm}YY}
            \hline
            \textbf{Dimension} & \textbf{L1 Regularisation} & \textbf{L2 Regularisation} \\
            \hline
            Mathematical penalty term    
            & \(\lambda \sum_{j=1}^{p} |w_j|\) 
            & \(\lambda \sum_{j=1}^{p} w_j^2\) \\
            Effect on model weights      
            & Drives irrelevant weights to zero (sparse) 
            & Shrinks irrelevant weights towards zero but rarely to zero (dense) \\
            Feature selection behaviour  
            & Performs feature selection by removing irrelevant features 
            & No feature selection. Retains all features \\
            Geometric interpretation     
            & Produces diamond-shaped constraint regions 
            & Produces circular constraint regions \\
            \hline
            \end{tabularx}
            \label{reg_comp}
            \end{table}

    \item   \begin{enumerate}[label=\roman*., topsep=0pt, itemsep=2pt]
                \item Advantage: L1 regularisation reduces irrelevant features (words) by setting their associated weights to zero. This creates a sparse data set which helps an SVM distinguish between
                      classes.

                      Disadvantage: L1 is rather unstable when features are correlated (similar words). L1 may select one feature and discards others, leading to loss of useful complimentary signals.
                      This can cause more variance in model performance across different training sets.
                    
                \item L2 regularisation is more appropriate when correlated features are present. Hence, when the goal is to achieve a robust model. L2 distributes weight across correlated features rather
                      than just choosing one. This creates a smoother decision boundary and helps the model generalise when many words have sentiment information that overlaps.
            \end{enumerate}
    
\end{enumerate}

\subsection{Question 2: Feature Extraction for Text Classification}

\subsection{Question 3: Support Vector Machines}